% Subsection 4.2.1: Prompt 模板与生成设置

基线采用单阶段 chat 格式 prompt，所有训练/推理组若涉及 Direct 推理侧均共享该模板，以保证训练侧主效应不被 prompt 差异污染。

\subsubsection{Chat 模板与系统指令}
基座为 Instruct 化模型，因此直接复用其官方 chat template：system 消息固化「\texttt{expert Python programmer}」角色定位，user 消息以 \texttt{Implement the function `\{func\_name\}'...} 句式承载题面与函数签名约束并以「Only output the Python code, no explanation.」收束，随后由 \texttt{apply\_chat\_template(add\_generation\_prompt=True)} 拼出最终输入序列。\texttt{func\_name} 与 \texttt{params} 从 3.1 节 \texttt{test\_info[0]} 的 \texttt{function\_name}、\texttt{parameter\_list} 字段直接取出，与评测时调用的被测函数同名同参，杜绝因函数名漂移导致的断言失败；逐字模板见附录~\ref{appendix:prompts}~A.2.1。

\subsubsection{生成采样参数}
基线的采样配置面向 pass@k 评测而非贪心解码，关键参数见表~\ref{tab:ch4-base-sampling}：温度取低值以兼顾多样性与正确率，每题采样次数共同决定了可估计到 pass@10 的最小采样规模；其余采样细节（核采样阈值、候选截断、停止符等）取常用默认值，详见仓库实现。

\begin{table}[!htb]
  \centering
  \caption{基线 Direct 推理的关键生成采样参数}
  \label{tab:ch4-base-sampling}
  \footnotesize
  \begin{tabular}{llp{6.8cm}}
    \toprule
    参数 & 取值 & 说明 \\
    \midrule
    每题采样数         & 10   & 同时支撑 pass@1/5/10 的无偏估计 \\
    温度（temperature） & 0.3  & 低温鼓励稳定输出，避免 pass@1 被随机性淹没 \\
    单次生成上限       & 1024 & 兼顾长解答与显存 \\
    \bottomrule
  \end{tabular}
\end{table}

上述参数在 4.3、4.4 节的所有 Direct 推理组（H4/H7/H9）中保持完全一致，确保在讨论训练侧效应时唯一变化的只有 LoRA 适配器本身。